% ---
% id: optics-2026-003
% title: Очень простая задача №3
% topic: observations
% subtopic: telescope-optics
% difficulty: 3
% source: camp
% event: Оптика I
% year: 2026
% grade: 9-10
% lang: ru
% ---
6 сентября 2022 года Юпитер оказался в «великом» (ближайшем к Солнцу
и Земле) противостоянии и подошёл к Земле на рекордно малое расстояние
за последние 59 лет — $\qty{591.3}{млн\ км}$. Это уникальный шанс для
получения наиболее детальных фотографий атмосферы планеты. Астроном-любитель
решил снять планету с помощью телескопа с диаметром объектива
$D = \qty{150}{мм}$ и относительным отверстием, равным $1:5$, используя
трёхкратную линзу Барлоу (ЛБ, увеличивающую эффективное фокусное расстояние
в 3 раза) и астрокамеру с размером пикселя матрицы
$\qty{2.4 \times 2.4}{мкм}$. Считать, что Юпитер обращается вокруг Солнца
в плоскости эклиптики. Определите:

\begin{enumerate}[label=\alph*)]
  \item масштаб (в угловых секундах на пиксель) и предел разрешающей
        способности оптической системы «телескоп–камера» (для длины волны
        $\lambda = \qty{555}{нм}$). Следует полагать, что величина
        атмосферных искажений во время наблюдения не меньше $\asec{1.5}$;
  \item максимальную допустимую выдержку для одного кадра видеоролика,
        на котором ещё не будет размытия изображения планеты в результате
        суточного движения Земли (при неподвижной трубе телескопа);
  \item оптимальные размеры одного кадра для съёмки Юпитера с использованием
        ЛБ из возможных форматов ($1936\times1096$, $1280\times960$,
        $640\times480$, $968\times548$, $320\times240$ пикселей), с целью
        минимизации размера файла;
  \item максимальную продолжительность видеоролика, из кадров которого ещё
        можно собрать изображение планеты без размытия деталей, обусловленного
        суточным вращением Юпитера. Следует полагать, что съёмка видеоролика
        осуществляется в режиме точного ведения трубы телескопа на
        роботизированной монтировке за Юпитером, компенсирующего суточное
        вращение Земли.
\end{enumerate}
